\section{Related Work}
\label{sec:related}

\para{Transformer regressors for puzzle difficulty.}
\glickformer~\citep{milosz2024glickformer} sets the current peer-reviewed state of the art for Lichess-puzzle rating regression at MAE 217.7 on a 42{,}000-puzzle held-out slice. It encodes each position along the solution as a 16-channel 8$\times$8 tensor and applies factorized spatio-temporal attention (following ViViT) with \chessformer~\citep{monroe2024chessformer} as the spatial backbone. Training uses uncertainty augmentation: labels are sampled from $\mathcal{N}(\text{Rating}, \text{RD}^2)$ clipped to $\pm3\text{RD}$. Follow-up work by the same group~\citep{milosz2024pretraining} explores pretraining objectives. Building on \glickformer's uncertainty-aware training recipe, we down-weight noisy-rating puzzles with sample weights $1/(1 + \text{RD}/40)$ but do not require the $10^7$-parameter transformer stack.

\para{Learned position/move sequence models.}
\ratingnet~\citep{omori2024ratingnet} uses a 4-layer CNN over piece-plane tensors followed by a BiLSTM over the move sequence to estimate player ratings from games, and is reported as ``competitive'' on the BigData Cup puzzle set. Ablations show that clock features contribute only $\approx$5\% of MAE reduction, motivating the interpretation that most of the signal is in position and move structure. Unlike \ratingnet, we use no learned position encoder: our board representation is a fixed 62-dim static feature vector, which we show is sufficient to reach within 23 \elo of the transformer state of the art when augmented with a lightweight tactical-density block.

\para{Skill-conditioned move models.}
\maia~\citep{mcilroy2020aligning} and Maia-2~\citep{jacob2024maia2} learn skill-conditioned move-prediction policies from human games at each Lichess Elo bracket. The third-place team in the IEEE BigData Cup~\citep{ieee2024cup} used the lowest \maia bracket at which the model's top move matches the puzzle solution as a difficulty proxy inside a \lightgbm regressor. We do not use \maia features in this iteration --- our tactical-density block is a cheaper CPU-only proxy for candidate-move ambiguity --- but note that \maia--bracket features are complementary and appear in our future-work list.

\para{Feature-engineered gradient boosting.}
The BigData Cup~\citep{ieee2024cup} report describes gradient-boosted trees with hand-crafted features and engine-derived signals (Stockfish evaluation drops, \maia bracket consistency) as the third-place approach; a pairwise learning-to-rank formulation~\citep{ieee2024pairwise} reframes the task as ranking rather than regression. In contrast, we position \ours as the first public, reproducible CPU-only gradient-boosting reference for puzzle-difficulty regression with a rich feature set (155 dims), and specifically isolate the contribution of tactical-density features via a controlled ablation against a 141-dim parent pipeline.

\para{Playing-strength chess transformers.}
\chessformer~\citep{monroe2024chessformer} introduces Smolgen, a learnable per-square attention bias that reflects chess movement patterns rather than Euclidean distance, and reaches 2347 Elo as a playing engine. Related work~\citep{zhang2024humanaligned} combines rating-conditioned move prediction with lightweight Stockfish rollouts. These works focus on strong play rather than difficulty prediction, but their intuition --- that piece-relationship structure carries more signal than raw square adjacency --- motivates our use of legal-move counts (which encode piece-relationship density) as difficulty features.

\para{Calibration for regression trees under label noise.}
Isotonic regression~\citep{barlow1972isotonic} and CDF quantile-matching are standard post-hoc calibrators for probabilistic and regression models. Squared-error tree ensembles under heavy label noise systematically compress tail predictions toward the mean; the CDF calibrator over-corrects and pays MAE. We use their \emph{convex blend} $\alpha \cdot \text{iso} + (1-\alpha)\cdot \text{cdf}$, whose blend weight is picked on a val-MAE-budgeted argmin over a mid-band calibration proxy. Both components are monotone, so Spearman rank order is preserved exactly.
